\subsubsection*{2.2.5.}
现在假设
\[
(*)\qquad\text{一般纤维 $X_\eta$ 光滑.}
\]
为简便起见, 还假设 $S$ 严格 Hensel. 利用 2.2.3 及 XII, 3.7 容易计算消灭圈层. 此处考虑的 $\Lambda$ 系数上同调, 可由所引段落中的 $\ell$ 进上同调通过万有系数公式(参见 2.1.13)得到. 结果如下.

\medskip
\noindent\underline{A. 无支集条件的上同调.}

对 $i\ne0,n$, 有 $R^i\psi_{\bar\eta}(\Lambda)=0$.

\begin{enumerate}
\item[1)] \underline{若 $n\ne0$:}
$\psi_{\bar\eta}(\Lambda)=\Lambda$. 此外, $R^n\psi_{\bar\eta}(\Lambda)$ (非典范地)同构于 $x_0$ 上的层 $\Lambda$ 在 $X_s$ 上的零延拓.

\item[2)] \underline{对任意 $n$:}
即使 $n=0$, 也可检验对 $i\ne n$ 有 $R^i\Phi_{\bar\eta}(\Lambda)=0$, 且 $R^n\Phi_{\bar\eta}(\Lambda)$ (非典范地)同构于 $x_0$ 上的层 $\Lambda$ 在 $X_s$ 上的零延拓.
\end{enumerate}

\medskip
\noindent\underline{B. 以 $x_0$ 为支集的上同调.}

有
\[
H^i_{\{x_0\}}\bigl(X_s,R\psi_{\bar\eta}(\Lambda)\bigr)=0
\qquad(i\ne n,2n).
\]
此外还有迹态射
\begin{equation}\tag{2.2.5.1}
\operatorname{tr}:
H^{2n}_{\{x_0\}}\bigl(X_s,R\psi_{\bar\eta}(\Lambda(n))\bigr)
\longrightarrow\Lambda.
\end{equation}
\underline{若 $n\ne0$:} 态射 (2.2.5.1) 是同构. 而且,
\[
H^n_{\{x_0\}}\bigl(X_s,R\psi_{\bar\eta}(\Lambda(n))\bigr)
\]
(非典范地)同构于 $\Lambda$.

\medskip
\noindent\underline{C. 对偶性.}
对
\[
a\in R^n\psi_{\bar\eta}(\Lambda)_{x_0},
\qquad
b\in H^n_{\{x_0\}}\bigl(X_s,R\psi_{\bar\eta}(\Lambda(n))\bigr),
\]
置 $(a,b)=\operatorname{Tr}(a\wedge b)$, 其中 $\operatorname{Tr}$ 为 XIII, 2.2.6 中的态射
\[
\operatorname{Tr}:
H^{2n}_{\{x_0\}}\bigl(X_s,R\psi_{\bar\eta}(\Lambda(n))\bigr)
\longrightarrow\Lambda.
\]
由 $X_{\bar\eta}$ 上的 Poincar\'e 对偶以及同构 2.2.3 可知, $(a,b)$ 给出自由 $\Lambda$-模
\[
R^n\psi_{\bar\eta}(\Lambda)_{x_0}
\quad\text{及}\quad
H^n_{\{x_0\}}\bigl(X_s,R\psi_{\bar\eta}(\Lambda(n))\bigr).
\]
之间的完美对偶.
\footnote{印本在所列的迹映射中漏去了 $\bar\eta$ 的上划线; 周围公式确定了此项的写法.}